\section*{Abstract}
\addcontentsline{toc}{section}{Abstract}

Particle-in-cell (PIC) simulations produce large datasets of particle trajectories that require automated methods for grouping and analysis. This internship adapts a previous spectrogram-based deep clustering workflow to six-channel records comprising particle position, momentum components, and local electromagnetic fields. The study uses 12,800 particles. After temporal subsampling and channel-wise standardization, the time series are converted into log-magnitude spectrograms. A convolutional autoencoder learns 32-dimensional representations, which are used for K-means initialization and deep clustering refinement with a reconstruction constraint. An elbow criterion selects five clusters from a candidate range of two to twelve. The method is compared with PCA followed by K-means on the same spectrogram inputs and a one-dimensional convolutional autoencoder with clustering refinement operating on standardized time series. On a common sample of 10,000 particles, the silhouette scores are 0.29957 for PCA + K-means, 0.79637 for the time-series method, and 0.87985 for the spectrogram method.  